%! Carrying Out a Test for a Population Mean — Unit 7, Lesson 7.5 — Homework
\documentclass[10pt]{article}
\usepackage{apstats-article}
\usepackage{apstats-boxes}
\usepackage{grading}
\gradingassignment{apstats-u07-l05-hw}
\gradingcourse{advanced_statistics}
\newcommand{\TallMath}[1]{$\displaystyle #1\rule[-1.4em]{0pt}{3.2em}$}

\begin{document}

\pageheader{Unit 7, Lesson 7.5}{Homework}

\vspace{0.15in}

\textit{Complete all four SPDC steps for each problem. Show all computations.}

\begin{enumerate}

\item \textbf{Plant Growth.} \quad
A botanist claims that a new fertilizer increases mean plant height beyond the historical
average of 12 cm after 4 weeks. She randomly assigns 25 plants to the new fertilizer.
$\bar x = 13.6$ cm, $s = 3.2$ cm. A dotplot of plant heights shows a roughly symmetric
distribution with no outliers.

\begin{enumerate}[label=\alph*.]
  \item \textbf{State:} $H_0$: \makebox[4cm]{\smash{\gradedblank[id=hw1a1, type=short, points=1, answer={{mu = 12}}, check=text, display={\boldmath$\mu = 12$}]{3.8cm}{0.19in}}} \quad $H_a$: \makebox[4cm]{\smash{\gradedblank[id=hw1a2, type=short, points=1, answer={{mu > 12}}, check=text, display={\boldmath$\mu > 12$}]{3.8cm}{0.19in}}}
  \item \textbf{Plan:} Method: \makebox[5cm]{\smash{\gradedblank[id=hw1b1, type=short, points=1, answer={{one-sample t-test for mu}}, check=text, display={\boldmath one-sample $t$-test for $\mu$}]{4.8cm}{0.19in}}}, $df = $ \makebox[1.5cm]{\smash{\gradedblank[id=hw1b2, type=short, points=1, answer=24, check=numeric, tolerance=0, display={\boldmath$24$}]{1.3cm}{0.19in}}}. Conditions:
        \begin{itemize}
          \boxguard
          \item Random: \hfill\makebox[4.6in][l]{\gradedhere{\gradedblank[id=hw1b3, type=rubric, points=1, display={random assignment of 25 plants --- met}]{4.6in}{0.19in}}}\\
                \gradedhide{\writelines{9}}
          \boxguard
          \item Normal ($n < 30$): \hfill\makebox[4.6in][l]{\gradedhere{\gradedblank[id=hw1b4, type=rubric, points=1, display={$n = 25 < 30$; dotplot symmetric, no outliers --- met}]{4.6in}{0.19in}}}\\
                \gradedhide{\writelines{9}}
        \end{itemize}
  \item \textbf{Do:} $SE = $ \makebox[3cm]{\smash{\gradedblank[id=hw1c1, type=short, points=1, answer=0.64, check=numeric, tolerance=0.005, display={\boldmath$0.64$}]{2.8cm}{0.19in}}} \quad $t \approx$ \makebox[2.5cm]{\smash{\gradedblank[id=hw1c2, type=short, points=1, answer=2.50, check=numeric, tolerance=0.02, display={\boldmath$2.50$}]{2.3cm}{0.19in}}} \quad
        $p \approx$ \texttt{tcdf}(\makebox[2cm]{\smash{\gradedblank[id=hw1c3, type=short, points=1, answer=2.50, check=numeric, tolerance=0.02, display={\boldmath$2.50$}]{1.8cm}{0.19in}}}, $10^{99}$, \makebox[2cm]{\smash{\gradedblank[id=hw1c4, type=short, points=1, answer=24, check=numeric, tolerance=0, display={\boldmath$24$}]{1.8cm}{0.19in}}}) $\approx$ \makebox[2cm]{\smash{\gradedblank[id=hw1c5, type=short, points=1, answer=0.010, check=numeric, tolerance=0.002, display={\boldmath$0.010$}]{1.8cm}{0.19in}}}
  \boxguard
  \item \textbf{Conclude:} \gradedhide{\writeline\\\gradedhere{\gradedblank[id=hw1d, type=rubric, points=2, display={$p \approx 0.010 < 0.05$: reject $H_0$. Convincing evidence the true mean height with the new fertilizer is $> 12$ cm.}]{\linewidth+6pt}{\gradedlinesheight{1}}}\writelines{9}}
\end{enumerate}

\vspace{0.08in}

\newpage
\item \textbf{Exam Score Study.} \quad
A university testing center claims the mean score on a standardized placement exam is 500.
A professor suspects the mean score at her campus is \emph{different} from 500. A random
sample of 45 students from her campus yields $\bar x = 485$, $s = 60$.

\begin{enumerate}[label=\alph*.]
  \item \textbf{State:} $H_0$: \makebox[4cm]{\smash{\gradedblank[id=hw2a1, type=short, points=1, answer={{mu = 500}}, check=text, display={\boldmath$\mu = 500$}]{3.8cm}{0.19in}}} \quad $H_a$: \makebox[4cm]{\smash{\gradedblank[id=hw2a2, type=short, points=1, answer={{mu != 500}}, check=text, display={\boldmath$\mu \ne 500$}]{3.8cm}{0.19in}}}
  \item \textbf{Plan:} Method: \makebox[5cm]{\smash{\gradedblank[id=hw2b1, type=short, points=1, answer={{one-sample t-test for mu}}, check=text, display={\boldmath one-sample $t$-test for $\mu$}]{4.8cm}{0.19in}}}, $df = $ \makebox[1.5cm]{\smash{\gradedblank[id=hw2b2, type=short, points=1, answer=44, check=numeric, tolerance=0, display={\boldmath$44$}]{1.3cm}{0.19in}}}. Conditions (brief):\par
        \gradedhide{\gradedhere{\gradedblank[id=hw2b3, type=rubric, points=1, display={random; $n \ge 30$}]{\linewidth+6pt}{0.19in}}\writelines{9}}
  \item \textbf{Do:} $SE = $ \makebox[3cm]{\smash{\gradedblank[id=hw2c1, type=short, points=1, answer=8.944, check=numeric, tolerance=0.01, display={\boldmath$8.944$}]{2.8cm}{0.19in}}} \quad $t \approx$ \makebox[2.5cm]{\smash{\gradedblank[id=hw2c2, type=short, points=1, answer=-1.68, check=numeric, tolerance=0.02, display={\boldmath$-1.68$}]{2.3cm}{0.19in}}}

        $p$ (two-sided) $\approx 2 \times \texttt{tcdf}($\makebox[3cm]{\smash{\gradedblank[id=hw2c3, type=short, points=1, answer=1.68, check=numeric, tolerance=0.02, display={\boldmath$1.68$}]{2.8cm}{0.19in}}}, $10^{99}$, \makebox[2cm]{\smash{\gradedblank[id=hw2c4, type=short, points=1, answer=44, check=numeric, tolerance=0, display={\boldmath$44$}]{1.8cm}{0.19in}}}$) \approx$ \makebox[2cm]{\smash{\gradedblank[id=hw2c5, type=short, points=1, answer=0.100, check=numeric, tolerance=0.005, display={\boldmath$0.100$}]{1.8cm}{0.19in}}}

  \boxguard
  \item \textbf{Conclude:} \gradedhide{\writeline\\\gradedhere{\gradedblank[id=hw2d, type=rubric, points=2, display={$p \approx 0.100 > 0.05$: fail to reject $H_0$. No convincing evidence the campus mean score differs from 500.}]{\linewidth+6pt}{\gradedlinesheight{1}}}\writelines{9}}
\end{enumerate}

\vspace{0.08in}

\newpage
\item \textbf{Matched Pairs: Study Techniques.} \quad
An educational researcher wants to know whether using flashcards improves quiz scores.
Ten students each take a quiz on two equivalent topics: once without flashcards and once
after using flashcards for one week. The differences $d_i = \text{flashcards} - \text{no flashcards}$
are: $\{4, 7, -1, 5, 9, 3, 6, 2, 8, 5\}$. \quad $\bar d = 4.8$, $s_d = 2.97$.
A boxplot of the differences is roughly symmetric with no outliers.

\begin{enumerate}[label=\alph*.]
  \item \textbf{State:} $\mu_d = $\par
        \gradedhide{\gradedhere{\gradedblank[id=hw3a1, type=rubric, points=1, display={true mean diff.\ (flash $-$ none)}]{\linewidth+6pt}{0.19in}}\writelines{9}}\par
        $H_0$: \makebox[3cm]{\smash{\gradedblank[id=hw3a2, type=short, points=1, answer={{mu_d = 0}}, check=text, display={\boldmath$\mu_d = 0$}]{2.8cm}{0.19in}}} \quad $H_a$: \makebox[3cm]{\smash{\gradedblank[id=hw3a3, type=short, points=1, answer={{mu_d > 0}}, check=text, display={\boldmath$\mu_d > 0$}]{2.8cm}{0.19in}}}
  \item \textbf{Plan:} Method: \makebox[5cm]{\smash{\gradedblank[id=hw3b1, type=short, points=1, answer={{matched-pairs t-test for mu_d}}, check=text, display={\boldmath matched-pairs $t$-test for $\mu_d$}]{4.8cm}{0.19in}}}, $df = $ \makebox[1.5cm]{\smash{\gradedblank[id=hw3b2, type=short, points=1, answer=9, check=numeric, tolerance=0, display={\boldmath$9$}]{1.3cm}{0.19in}}}. Conditions:
        \begin{itemize}
          \boxguard
          \item Independence: \hfill\makebox[4.6in][l]{\gradedhere{\gradedblank[id=hw3b3, type=rubric, points=1, display={students tested independently; differences independent --- met}]{4.6in}{0.19in}}}\\
                \gradedhide{\writelines{9}}
          \boxguard
          \item Normal ($n < 30$): \hfill\makebox[4.6in][l]{\gradedhere{\gradedblank[id=hw3b4, type=rubric, points=1, display={$n = 10 < 30$; boxplot symmetric, no outliers --- met}]{4.6in}{0.19in}}}\\
                \gradedhide{\writelines{9}}
        \end{itemize}
  \item \textbf{Do:} $SE_d = \dfrac{s_d}{\sqrt{n}} \approx$ \makebox[2.5cm]{\smash{\gradedblank[id=hw3c1, type=short, points=1, answer=0.939, check=numeric, tolerance=0.005, display={\boldmath$0.939$}]{2.3cm}{0.19in}}} \quad
        $t = \dfrac{\bar d - 0}{SE_d} \approx$ \makebox[2.5cm]{\smash{\gradedblank[id=hw3c2, type=short, points=1, answer=5.11, check=numeric, tolerance=0.03, display={\boldmath$5.11$}]{2.3cm}{0.19in}}}

        $p \approx$ \makebox[4cm]{\smash{\gradedblank[id=hw3c3, type=short, points=1, answer=0.0003, check=numeric, tolerance=0.0002, display={\boldmath$0.0003$}]{3.8cm}{0.19in}}}

  \boxguard
  \item \textbf{Conclude:} \gradedhide{\writeline\\\gradedhere{\gradedblank[id=hw3d, type=rubric, points=2, display={$p \approx 0.0003 < 0.05$: reject $H_0$. Convincing evidence flashcards raise mean quiz score ($\mu_d > 0$).}]{\linewidth+6pt}{\gradedlinesheight{1}}}\writelines{9}}
\end{enumerate}

\end{enumerate}

\begin{reflectionbox}
A student says: ``We got $p = 0.08$, so we know the null hypothesis is probably true.''
Explain what is wrong with this statement. What does ``fail to reject $H_0$'' actually mean?
\writelines{3}
\end{reflectionbox}

\begin{extensionbox}
For Problem 2, suppose the professor changes her hypotheses to $H_a: \mu < 500$ after seeing
that $\bar x = 485$. (a) Explain why this is a problem. (b) Compute the new (incorrect)
$p$-value using the one-sided left tail. (c) Compare it to the two-sided result.
\writelines{3}
\end{extensionbox}

\end{document}
